% SHARD-ID: CA-86-SPACETIME-TROTTER-DICTIONARY
% SHARD-TITLE: String-Net Histories as Trotter Expansions: the Spacetime Dictionary
% SHARD-SUMMARY: Proves that the brick-wall Trotter expansion of an anyonic chain Hamiltonian built from morphisms of C is a sum over string-net histories whose amplitude factorises into a C-evaluation obeying the Levin-Wen rules times a product of tile weights that breaks isotopy.
% SHARD-SUMMARY: Calibrates the dictionary on the dense and dilute Temperley-Lieb integrable lines, where the tile weights are the sourced face weights and the Hamiltonian limit, the isotropic point and Nienhuis's honeycomb critical point lie in one commuting family.
% SHARD-SUMMARY: Derives x = 1/(2 cos lambda) at u = lambda, reproducing the Nienhuis critical fugacity on branch 1 and the second-branch value on branch 2, and states the Fibonacci net-gas frontier beyond loops.
% SHARD-KEYWORDS: Levin-Wen, string-net, Trotter, transfer matrix, spacetime history, dilute O(n), Nienhuis, honeycomb, integrable line, Fibonacci, net gas

\section{String-Net Histories as Trotter Expansions: the Spacetime Dictionary}
\label{sec:spacetime-trotter-dictionary}

\subsection{Status, Sources, and Dependencies}

\textbf{Status: derived.}  Lemma 86.1 and
Proposition 86.4 are checked local derivations; the face
weights, branch data and central charges they use are source-local; the
final subsection is a proposal.  Reproducibility: the trigonometric identities
of Proposition 86.4 were evaluated numerically for
\(L=2,\dots,5\) on both branches (worklog chunk 016); a Julia port is an
open item.

Dependencies: CA-65, CA-66, CA-69, CA-74; CONVENTIONS.md (a), (b), (r), (ac).

% Source: references/string-net/LevinWen2005/source/strnet.tex:497--524 -- the
%   four local rules; :548--576 -- their RG motivation; :777--791, :806--813 --
%   a_s = d_s / sum d_i^2 is the unique phase with a smooth continuum limit.
% Source: references/lattice-symmetry/KooSaleur1993/source/9312156.tex:318--323,
%   :724--738 -- sqrt(Q) = 2 cos gamma; log of the diagonal TL transfer matrix
%   = (u/sin gamma) sum_j e_j + O(u^2), a commuting family.
% Source: references/lattice-symmetry/ZhouBatchelor1997/source/9611156.tex:105--145 --
%   dilute O(n) loop weights rho_1..rho_9(u), n = -2 cos 4 lambda; honeycomb
%   Nienhuis model at u = lambda, 2 lambda; :237--262 -- dilute A_L face
%   weights, n = 2 cos(pi/(L+1)); :264--276 -- the four branches; :388--399,
%   :430--455 -- c = 1 - 6/(h(h-1)), h = L+2 (branch 1), L+1 (branch 2).
% Source: references/lattice-symmetry/DiluteA22Strip2022/source/stripA22.tex:669--745 --
%   face operator = nine tiles in dTL_2(beta); :800--806, :849 -- crossing
%   symmetry u <-> 3 lambda - u; identity-arc factorisation at u = 0.
% Source: references/lattice-symmetry/ChelkakGlazmanSmirnov2016/source/StressTensor.tex:561,
%   :572--584, :600 -- Nienhuis x_c(n) = 1/sqrt(2 + sqrt(2-n)), second branch
%   1/sqrt(2 - sqrt(2-n)), c(kappa) at x_c and x > x_c.
% Source: references/text/GoldenChainFeiguinEtAl.txt:21--26, :238--242 --
%   golden chain critical, c = 7/10, exact map to the k = 3 RSOS model.
% Source: references/cft/AasenFendleyMong2020/source/_Statistical_mechanics_from_fusion_categories.tex:15--20,
%   :537--553 -- weights sum_chi A(chi) P^(chi); Hamiltonian from A = 1 - eps A_j;
%   beyond Temperley-Lieb, integrability requires further fine-tuning.

\subsection{The Levin--Wen Rules and What They Fix}

\textbf{Source-local.}  Levin and Wen specify a fixed-point wave function
\(\Phi\) on string-net configurations by four local linear relations:
(1) topological invariance, (2) a closed type-\(i\) loop contributes
\(d_i\), (3) a bubble on a string is \(\delta_{ij}\), (4) the \(F\)-move.
Rules (2)--(3) are motivated by scale invariance of an RG fixed point, rule
(1) by the expectation that topological phases have topologically invariant
fixed points, and rule (4) is the simplest completing constraint.  The
plaquette coefficients \(a_s=d_s/\sum_i d_i^2\) are not chosen: they are the
unique choice with a smooth continuum limit.  The lesson carried forward is
that a \emph{principle} (scale invariance of the fixed point) fixes the model.

\subsection{The Trotter Dictionary}

\textbf{Setting (CONVENTIONS (r), (ac)).}  \(\Cc\) unitary fusion, \(O=\one\oplus X\), \(A_L=\operatorname{End}_{\Cc}(O^{\otimes L})\).
A nearest-neighbour Hamiltonian is \(H=-\sum_{j}h_j\) with
\(h_j\in\operatorname{End}_{\Cc}(O\otimes O)\) acting on sites \(j,j+1\).
Fix a basis of \emph{elementary tiles} \(D_\alpha\) of
\(\operatorname{End}_{\Cc}(O\otimes O)\) consisting of morphism diagrams
(vacancy projectors, through-strings, cup, cap, hops, and for non-loop
categories the fusion vertices), and expand \(h_j=\sum_\alpha c_\alpha D_\alpha^{(j)}\)
with \(c_\alpha\in\mathbb{C}\).  Brick-wall Trotter with step \(\delta\tau\):
\[
  U(\delta\tau)=\prod_{j\ \mathrm{odd}}e^{\delta\tau h_j}\prod_{j\ \mathrm{even}}e^{\delta\tau h_j},
  \qquad
  e^{-\beta H}=U(\delta\tau)^{M}+O(\beta\,\delta\tau)\ \ (\beta=M\delta\tau)
\]
on a finite chain, by the Lie--Trotter formula.  A \emph{history} \(\Gamma\)
assigns to each plaquette of the \(M\)-row brick lattice either the identity
tile or one elementary tile \(D_{\alpha(p)}\).  Composing the tiles row by
row gives a morphism \(\operatorname{Ev}_{\Cc}(\Gamma)\in A_L\), the
\(\Cc\)-evaluation of the spacetime string-net drawn by \(\Gamma\).

\textbf{Lemma 86.1 [derived; tile expansion].}  To first order per plaquette,
\[
  U(\delta\tau)^{M}
  =\sum_{\Gamma}\Big(\prod_{p:\,\alpha(p)\neq\mathrm{id}}\delta\tau\,c_{\alpha(p)}\Big)\,
  \operatorname{Ev}_{\Cc}(\Gamma)\;+\;R,
\]
where \(R\) collects every term carrying a second-order factor
\(O(\delta\tau^2\|h_j\|^2)\) on at least one plaquette.
\textbf{Proof.}  \(e^{\delta\tau h_j}=1+\delta\tau h_j+O(\delta\tau^2\|h_j\|^2)\) per plaquette.
Expanding \(\prod_p(1+\delta\tau h_{j(p)})\) and each \(h_j\) in tiles gives
exactly the sum over histories with one tile per plaquette; composition of
the tiles in the brick order is the composition in \(A_L\), which is
\(\operatorname{Ev}_{\Cc}(\Gamma)\).  The remainder collects all terms with a
second-order factor on at least one plaquette.  \(\square\)

\textbf{Corollary 86.2 [the Levin--Wen rules are automatic].}  \(\operatorname{Ev}_{\Cc}(\Gamma)\) obeys rules (2), (3), (4): a closed
\(X\)-loop evaluates to \(d_X\) by the raw cup normalisation
\(v^\dagger v=d_X\operatorname{id}_\one\) of (r); a bubble between simple
labels is \(\delta\) by Schur; recoupling is the \(F\)-move.  Rule (1),
isotopy, holds for \(\operatorname{Ev}_{\Cc}\) alone (planar isotopy of
morphism diagrams in a pivotal category) and fails for the full amplitude,
which carries the geometric factor \(\prod_p\delta\tau\,c_{\alpha(p)}\)
depending on how many plaquettes each string traverses.  Hence
\[
  \text{amplitude}(\Gamma)=\underbrace{\operatorname{Ev}_{\Cc}(\Gamma)}_{\text{topological}}
  \times\underbrace{\textstyle\prod_p w(p)}_{\text{geometric}},
  \qquad w(\mathrm{id})=1,\ w(\alpha)=\delta\tau\,c_\alpha .
\]
This is the spacetime reading made precise: rules (2)--(4) intact, rule (1)
replaced by anisotropic tile weights (identity tiles weigh 1, all others
\(O(\delta\tau)\)); isotropy can only emerge in the continuum.

\textbf{Proposition 86.3 [derived; exactness on an integrable line].}  Suppose the tiles carry weights \(\rho_\alpha(u)\) forming a face operator
\(X_j(u)=\sum_\alpha\rho_\alpha(u)D^{(j)}_\alpha\) with \(X_j(0)=N(0)\cdot 1\),
whose diagonal-to-diagonal (brick-wall) transfer matrices
\(T(u)=\prod_{j\ \mathrm{odd}}X_j(u)\prod_{j\ \mathrm{even}}X_j(u)\) commute for
all \(u\), as the sourced \(\hat\tau_D\) does.  Then
\(h_j=X_j'(0)/N(0)-(N'(0)/N(0))\,1\) satisfies
\(1+\delta\tau h_j=X_j(\delta\tau)/N(\delta\tau)+O(\delta\tau^2)\), the
Trotter product is \(T(\delta\tau)^M\) up to normalisation, and \(H\)
commutes with \(T(u)\) for every \(u\).  In particular the ground state of
\(H\) is a common eigenvector of the transfer matrices at every \(u\),
including the isotropic point.

\textbf{Proof.}  The first two claims are Taylor expansion.  For the third, the Taylor
coefficients of a commuting one-parameter family commute with every member
of the family, and \(H\) is the first coefficient of \(-\log T(u)\).  \(\square\)
On an integrable line no perturbative matching of rates is needed: \(u\)
interpolates exactly between the Hamiltonian limit and the isotropic point.

\subsection{Dense Calibration}

Take \(X\) with \(d=d_X\) and \(H=-\sum_j e_j\), \(e_j=b_jd_j\) the cup-cap
generator of CA-66 (\(e_j^2=d\,e_j\)).  Because \(e_j^2\propto e_j\),
\[
  e^{\beta e_j}=1+w(\beta)\,e_j,\qquad w(\beta)=\frac{e^{\beta d}-1}{d}=\beta+O(\beta^2),
\]
so the brick-wall expansion is \emph{exact} with two tiles per plaquette and
\[
  \langle\mathrm{out}|U^{M}|\mathrm{in}\rangle
  =\sum_{\Gamma}w^{\#\{e\text{-tiles}\}}\,d^{\#\{\text{closed loops}\}}\,
  \langle\mathrm{out}|\Gamma_{\mathrm{open}}|\mathrm{in}\rangle .
\]
This is the sourced diagonal Temperley--Lieb transfer matrix at small
spectral parameter: with \(d=\sqrt{Q}=2\cos\gamma\),
\(\log(Q^{-L/2}\epsilon_y^L\hat\tau_D(u))=(u/\sin\gamma)\sum_je_j+O(u^2)\),
a commuting family, so Proposition 86.3 applies.  For
\(d=\varphi\), \(\gamma=\pi/5\), the antiferromagnetic sign is the golden
chain, critical with \(c=7/10\) by the exact RSOS map.  The dense chain has
no continuous parameter at all: the category fixes \(\gamma\) and the
remaining choice is a sign.

\subsection{Dilute Calibration}

\textbf{Tiles (CONVENTIONS (ac)).}  The dilute face operator is a linear
combination of nine tiles.  With in-nodes on the left and bottom edges
(sites \(j\), \(j+1\)) and out-nodes on the top and right edges, the tiles
are, in the source's numbering:
\(1\): both vacant, \(P_{00}\);
\(2,3\): one through-string, \(P_{X0},P_{0X}\);
\(8\): two through-strings, \(P_{XX}\);
\(9\): cap-then-cup, \(e_j\);
\(4,5\): cap \(d_j\) and cup \(b_j\);
\(6,7\): hops \(h_j^{\pm}\).
The sourced weights, normalised so that \(\rho_{1,2,3,8}(0)=1\) and
\(\rho_{4,\dots,7,9}(0)=0\), are
\[
\begin{aligned}
\rho_1&=1+\tfrac{\sin u\,\sin(3\lambda-u)}{\sin2\lambda\sin3\lambda},&
\rho_{2,3}&=\tfrac{\sin(3\lambda-u)}{\sin3\lambda},&
\rho_{4,5}&=\tfrac{\sin u}{\sin3\lambda},\\
\rho_{6,7}&=\tfrac{\sin u\,\sin(3\lambda-u)}{\sin2\lambda\sin3\lambda},&
\rho_8&=\tfrac{\sin(2\lambda-u)\sin(3\lambda-u)}{\sin2\lambda\sin3\lambda},&
\rho_9&=-\tfrac{\sin u\,\sin(\lambda-u)}{\sin2\lambda\sin3\lambda},
\end{aligned}
\]
with loop fugacity \(n=-2\cos4\lambda\).  The \(u=0\) values confirm the
tile assignment: \(X_j(0)=P_{00}+P_{X0}+P_{0X}+P_{XX}=1\), matching the
sourced factorisation into identity-arc triangles.  The Hamiltonian limit
\(H_{\mathrm{dil}}=-\sum_jX_j'(0)\) is
\[
  H_{\mathrm{dil}}=\sum_j\Big[
  -\tfrac{P_{00}}{\sin2\lambda}
  +\cot3\lambda\,(P_{X0}+P_{0X})
  +\tfrac{\sin5\lambda}{\sin2\lambda\sin3\lambda}P_{XX}
  -\tfrac{b_j+d_j}{\sin3\lambda}
  -\tfrac{h_j^++h_j^-}{\sin2\lambda}
  +\tfrac{\sin\lambda}{\sin2\lambda\sin3\lambda}\,e_j\Big].
\]
Every coupling is a function of \(\lambda\) alone.  For a Temperley--Lieb
type category with \(d=2\cos(\pi/(L+1))\) the equation \(n=d\) has the two
solutions \(\lambda=\tfrac{\pi}{4}\tfrac{L}{L+1}\) (branch 1) and
\(\lambda=\tfrac{\pi}{4}\tfrac{L+2}{L+1}\) (branch 2), both critical, with
\(c=1-6/(h(h-1))\), \(h=L+2\) resp.\ \(h=L+1\).  Fibonacci is \(L=4\):
branch 1 gives \(c=4/5\), branch 2 gives \(c=7/10\).

\textbf{Proposition 86.4 [derived; honeycomb reduction at \(u=\lambda\)].}  At \(u=\lambda\) the nine-tile model is the honeycomb \(O(n)\) loop model
with edge fugacity \(x=1/(2\cos\lambda)\).  On branch 1 this is Nienhuis's
critical point \(x_c(n)=1/\sqrt{2+\sqrt{2-n}}\); on branch 2 it is the
second-branch value \(1/\sqrt{2-\sqrt{2-n}}\).

\textbf{Proof.}  Split each square tile along the diagonal from the bottom-left to the
top-right corner into two trivalent vertices, \(A\) with edges
\(\{\mathrm{top},\mathrm{left},\mathrm{diag}\}\) and \(B\) with
\(\{\mathrm{bottom},\mathrm{right},\mathrm{diag}\}\).  Give each occupied
half-edge of a tile weight \(x^{1/2}\) and the internal diagonal weight
\(x\).  Then the tiles weigh: \(1\mapsto1\); \(2,3\mapsto x\) (an arc inside
\(A\) or \(B\)); \(4,5,6,7\mapsto x^2\) (an arc through the diagonal);
\(8\mapsto x^2\); \(9\mapsto0\) (it would use the diagonal twice).  At
\(u=\lambda\): \(\rho_9=0\); \(\rho_4=\rho_5=\rho_6=\rho_7=\rho_8=\sin\lambda/\sin3\lambda\);
\(\rho_{2,3}=\sin2\lambda/\sin3\lambda\);
\(\rho_1=1+\sin\lambda/\sin3\lambda=2\sin2\lambda\cos\lambda/\sin3\lambda\).
Hence \(\rho_2/\rho_1=1/(2\cos\lambda)=:x\) and
\(\rho_4/\rho_1=\sin\lambda/(4\sin\lambda\cos^2\lambda)=x^2\), consistent.
Finally \(x^{-2}=4\cos^2\lambda=2+2\cos2\lambda\) and
\(\sqrt{2-n}=\sqrt{2+2\cos4\lambda}=2|\cos2\lambda|\), so \(x=x_c\) when
\(\cos2\lambda\ge0\) (branch 1, \(\lambda<\pi/4\)) and \(x=1/\sqrt{2-\sqrt{2-n}}\)
when \(\cos2\lambda<0\) (branch 2).  \(\square\)

\textbf{Checked arithmetic (Fibonacci, \(L=4\)).}
\[
\begin{array}{c|c|c|c|c|c}
\text{branch} & \lambda & x=1/(2\cos\lambda) & \text{Nienhuis value} & c & \text{couplings }(-\rho_1',-\rho_2',-\rho_4',-\rho_6',-\rho_8',-\rho_9')\\\hline
1 & \pi/5 & 0.618034 & x_c=0.618034 & 4/5 & (-1.0515,\ -0.3249,\ -1.0515,\ -1.0515,\ 0,\ +0.6498)\\
2 & 3\pi/10 & 0.850651 & \tilde x_c=0.850651 & 7/10 & (-1.0515,\ -3.0777,\ -3.2361,\ -1.0515,\ -3.4026,\ +2.7528)
\end{array}
\]
The same identities hold for \(L=2,3,5\) (worklog 016).  Note the sign of
the \(e_j\) term: the integrable dilute chain has a \emph{repulsive}
cup-cap term on both branches, opposite to the dense golden chain, and on
branch 1 the \(P_{XX}\) coupling vanishes because \(\sin5\lambda=\sin\pi=0\).

\subsection{What the Calibration Says About the Principle}

For Temperley--Lieb type input the machine closes with no continuous
choice.  The category fixes \(n=d\), hence \(\lambda\) up to a branch; the
couplings of \(H_{\mathrm{dil}}\) are functions of \(\lambda\); the
Hamiltonian, the isotropic point \(u=3\lambda/2\) and the honeycomb point
\(u=\lambda\) are one commuting family; the limit is critical with a sourced
\(c\).  Isotropy is restored along the family, not imposed.  The geometric
parameter is the anisotropy \(u\), not a coupling.  The residual freedom is
the branch: two critical points for \(u>0\), two more with an Ising factor
for \(u<0\).  Whether a principle selects among branches is open.

\subsection{Beyond Loops: the Fibonacci Net Gas (Proposal)}

For Fibonacci on \(O=\one\oplus\tau\), \(O\otimes O=2\cdot\one\oplus3\cdot\tau\), so
\(\operatorname{End}(O\otimes O)=M_2\oplus M_3\) has dimension 13, while
\(\mathrm{dTL}_2(\varphi)\) has dimension 9 (CA-69).  The four missing
dimensions are the vertex tiles \(\tau\otimes\tau\leftrightarrow\tau\) with
one vacancy, which CA-74 proves lie outside the dilute loop algebra.  The
general reflection-symmetric nearest-neighbour Hamiltonian therefore has,
beyond the loop couplings above, a vertex coupling \(v\):
\[
  h=\big[\text{diag}(\epsilon_{00},\epsilon_{1})+w(b+d)\big]_{\one\text{-block}}
  \oplus\big[\text{diag}(\mu,\mu,\epsilon_{\tau})+t(h^++h^-)+v(\nu+\nu^\dagger)\big]_{\tau\text{-block}},
\]
with \(\nu:\ket{\tau\tau\to\tau}\mapsto\ket{\tau\one}+\ket{\one\tau}\) built
from the vertex morphism of \(\Cc\) (normalisation to be fixed in
CONVENTIONS before use).  Lemma 86.1 turns any such \(h\) into a
\emph{trivalent net gas} in spacetime with per-tile weights
\(\delta\tau\cdot(\text{coupling})\) and the \(\Cc\)-evaluation of the net.
Two corners are calibrated on the quantum side: \(\mu\to-\infty\) freezes
every site to \(\tau\) and recovers the dense chain (\(c=7/10\) or \(4/5\)
by sign); the \(v=0\) hyperplane is the parity-preserving dilute loop
sector of CA-74, which contains \(H_{\mathrm{dil}}\).  No integrable line is known once \(v\neq0\): the sourced statement is
that beyond Temperley--Lieb, integrability requires further fine-tuning.
Proposition 86.3 is therefore unavailable, and the
comparison of rates is genuinely first-order: the open problem is to find
the locus in \((\epsilon,\mu,w,t,v)\) at which the net gas is critical.  The
first analytical tool to point at it is a lattice holomorphicity condition
on a \(Z(\Cc)\)-spin defect endpoint, which is the next shard.
